We sweep one channel dimension at a time for the Gumbel-softmax system (3 seeds, same budget). Table~\ref{tab:sv} and Fig.~\ref{fig:sv} vary $V$ at $L=5$; Table~\ref{tab:sl} and Fig.~\ref{fig:sl} vary $L$ at $V=6$.

\begin{table}[t]\centering\caption{Vocabulary sweep, $L=5$.}\label{tab:sv}
\resizebox{\columnwidth}{!}{\begin{tabular}{lcccc}
\toprule
Task & held-out & train & topsim & unique msgs \\
\midrule
V2\_L5 & 0.000 $\pm$ 0.000 & 0.023 $\pm$ 0.014 & 0.264 $\pm$ 0.088 & 5.333 $\pm$ 3.215 \\
V3\_L5 & 0.000 $\pm$ 0.000 & 0.030 $\pm$ 0.017 & 0.136 $\pm$ 0.047 & 8.667 $\pm$ 4.933 \\
V4\_L5 & 0.000 $\pm$ 0.000 & 0.126 $\pm$ 0.079 & 0.179 $\pm$ 0.028 & 40.333 $\pm$ 20.257 \\
V8\_L5 & 0.064 $\pm$ 0.059 & 0.474 $\pm$ 0.363 & 0.236 $\pm$ 0.161 & 125.667 $\pm$ 94.817 \\
V16\_L5 & 0.487 $\pm$ 0.247 & 0.949 $\pm$ 0.021 & 0.388 $\pm$ 0.032 & 242.000 $\pm$ 6.245 \\
\bottomrule
\end{tabular}
}\end{table}
\begin{table}[t]\centering\caption{Length sweep, $V=6$.}\label{tab:sl}
\resizebox{\columnwidth}{!}{\begin{tabular}{lcccc}
\toprule
Task & held-out & train & topsim & unique msgs \\
\midrule
V6\_L2 & 0.013 $\pm$ 0.022 & 0.038 $\pm$ 0.003 & 0.122 $\pm$ 0.084 & 9.667 $\pm$ 1.155 \\
V6\_L3 & 0.013 $\pm$ 0.022 & 0.094 $\pm$ 0.073 & 0.152 $\pm$ 0.091 & 29.333 $\pm$ 19.858 \\
V6\_L4 & 0.000 $\pm$ 0.000 & 0.172 $\pm$ 0.064 & 0.203 $\pm$ 0.072 & 51.667 $\pm$ 19.858 \\
V6\_L6 & 0.141 $\pm$ 0.118 & 0.690 $\pm$ 0.072 & 0.334 $\pm$ 0.018 & 187.000 $\pm$ 15.133 \\
V6\_L8 & 0.077 $\pm$ 0.102 & 0.420 $\pm$ 0.225 & 0.245 $\pm$ 0.026 & 126.000 $\pm$ 50.110 \\
\bottomrule
\end{tabular}
}\end{table}

\begin{figure}[t]\centering\includegraphics[width=0.9\columnwidth]{sweep_vocab_all.pdf}
\caption{Train and held-out accuracy and topsim versus vocabulary size ($L=5$).}\label{fig:sv}\end{figure}
\begin{figure}[t]\centering\includegraphics[width=0.9\columnwidth]{sweep_length_all.pdf}
\caption{Train and held-out accuracy and topsim versus message length ($V=6$).}\label{fig:sl}\end{figure}

Larger vocabulary monotonically raises train accuracy (0.023 at V2\_L5 to 0.949 at V16\_L5) and held-out accuracy (0.000 to 0.487), and topsim increases from 0.136 at V3\_L5 to 0.388 at V16\_L5 (non-monotonic at V2\_L5, 0.264, where only a handful of distinct messages are used). The length sweep is non-monotonic: V6\_L6 has the highest train accuracy (0.690) and V6\_L8 is lower (0.420); their held-out accuracies (0.141 vs 0.077) are within one standard deviation at $n=3$. Longer messages may be harder to optimise rather than limited by capacity, but we did not test this; the V4\_L4 control (Sec.~\ref{sec:results}) shows that longer training alone does not rescue the Gumbel-softmax system there. Stronger channels also use more distinct messages (n\_unique\_messages in the tables), i.e.\ they drift towards one message per object, which is what the random code does by construction. The best topsim anywhere for a learned system is 0.388, far from the oracle's 1.000.
